\subsection{Eleven positive unknowns: an auxiliary offset and two root centers lower degree to 211}
\label{sec:positive-eleven-offset}

The following construction has eleven strictly positive integer unknowns,
exact degree $211$, and $112$ live binary arithmetic operations. The
centered predecessor in the incoming joint-guards package has degree $257$
in $114$ operations; the uncentered Pell-descent version has degree
$277$ in $111$ operations. The improvement is therefore a strict improvement in both degree and evaluator size over the centered predecessor.

The new step is the affine substitution
\[
 T_{\rm old}=T+F+\alpha.
\]
It makes the auxiliary quantity positive at an earlier stage of the proof,
so an expensive exceptional-sign term can be removed from the transport
weight. A second center absorbs the strong root, leaving a degree-two
gap bound. This is an existence-preserving substitution on canonical source
witnesses, not a bijection of the entire positive zero sets.

\subsection{The complete polynomial}

Use the coherent compact44 fixed tuple $(\beta,K_0,M_C,M_F)$, where
$\beta=2^{590069}-1$ and $M_F$ is already shifted. The tuple is defined by
the retained finite recipe and actual integer payloads. The unknowns are
\[
 j,F,\alpha,i,T,X,Y,n,\eta,\zeta,\gamma>0.
\]
Every remaining symbol is a polynomial abbreviation. Define
\begin{align*}
 J&=3j+1,\qquad q=\beta J+1,\qquad C=q-F-\alpha,\\
 E&=XY,\qquad k=\eta+\zeta,\qquad c=kY+\eta,\\
 a&=E+Y,\qquad H=4a+3,\qquad\Delta=a^2+H=(a+2)^2-1,\\
 \ell&=EkY,\qquad I=ic^2,\qquad M=q^3(q-1),\\
 R&=((q-1)(q-F)+\alpha)(q^2-1)+(M_C+qM_F)J,\\
 B_t&=(qK_0+X)C-qF,\qquad D_\tau=\ell(\ell+k)+1,\\
 D_f&=1+\Delta I^2,\qquad Q=\Delta^2I^2,\qquad\delta=Q-1,\\
 T_{\rm old}&=T+F+\alpha,\qquad C_f=(a+2)I,\\
 V&=D_f(cT_{\rm old}-R)-c,\qquad L=kT_{\rm old},\\
 W&=1+(q-1)L,\\
 \overline Z&=M(\ell+k+n+L+V+C_f)+(q-1)Y-q^3WB_t,\\
 r_0&=k-R-1,\qquad S=4qY,\qquad U=qS,\\
 D_*&=U(X+ac)+H(U\gamma+SX-r_0),\\
 A_1&=D_*^2-U^2(1+\Delta c^2).
\end{align*}
The three-root norm is
\begin{align*}
 A&=M^2D_\tau,\qquad B_f=M^2D_f,\qquad C_{\rm num}=M^2(QV^2-1),\\
 T_N&=\delta^2AB_f,\qquad R_N=\delta C_{\rm num},\\
 J_N&=\delta(\overline Z^2-A-B_f)-C_{\rm num},\\
 N_c&=J_N^2-4[T_N+R_N(A+B_f)],\\
 N_8&=N_c^2-64T_NR_N\overline Z^2.
\end{align*}
The proposed single polynomial equation is
\begin{equation}\label{eq:offset-p211}
 \boxed{P_{11}=N_8+A_1+N_8A_1=0.}
\end{equation}
Its literal graph is \texttt{p11\_positive\_degree211.json}.

\subsection{Recovering radicals before any transport sign}

Put, only in this proof,
\[
 t=\sqrt{D_\tau},\qquad f=\sqrt{D_f},\qquad
 y=\sqrt{\frac{QV^2-1}{Q-1}}.
\]
The unconditional three-root separation argument of source70 continues to
apply: it requires only $T_{\rm old}>0$ and $I\ge c^2$, not positivity of
$C$. We recall its substance. Let $b=cT_{\rm old}-R$. If $b=0$, then
$V=-c$, $y<c+1$, and $f-t>(a+1)c(c-1)>c+1$. If $b\ne0$, then
\[
 y>|V|\ge D_f-c>2f>f+t.
\]
Here $f>3c^2$ and $f>t$ follow directly from $\Delta\ge15$ and
$c>kY$. Thus one of the three roots strictly dominates the other two.
Their positive rational radicands are pairwise distinct: $D_\tau<D_f<Q$,
and equality of the third radicand to either integer $D\in\{D_\tau,D_f\}$
would imply $Q(D-V^2)=D-1$, impossible since $0<D-1<Q$.

A rational signed sum of three such roots forces them all rational. For
signed roots $x_1,x_2,x_3$ with rational nonzero sum $s$ and squares $a,b,c$,
set
\[
 u=\frac{s^2-a-b-c}{2},\qquad
 v=\frac{u^2-ab-ac-bc}{2s}.
\]
Then
\[
 x_1=\frac{a(a+u)}{as+v},\qquad
 as+v=x_1(x_1+x_2)(x_1+x_3)\ne0,
\]
and the cyclic formulas prove the assertion.

The universal norm identity
\[
 \bigl[(x^2-a-b-c)^2-4(ab+ac+bc)\bigr]^2-64abcx^2
 =\prod_{\varepsilon\in\{-1,1\}^3}
 (x-\varepsilon_1\sqrt a-\varepsilon_2\sqrt b-\varepsilon_3\sqrt c)
\]
therefore shows that $N_8=0$ supplies rational signed roots with
\[
 \frac{\overline Z}{M}=\tau_*+f_*+y_*,\qquad
 \tau_*^2=D_\tau,\quad f_*^2=D_f,\quad
 QV^2-(Q-1)y_*^2=1.
\]
The first two roots are integers. The adjacent-coefficient Pell descent
of source69 makes $y_*$ an integer as well. In that descent, writing
$y_*=a/m$ in lowest terms gives $m^2\mid Q-1$. After replacing $V$ and
$y_*$ by their absolute values at each step, the transformation
\[
 (V,y_*)\longmapsto
 ((2Q-1)V-2(Q-1)y_*,\ (2Q-1)y_*-2QV)
\]
preserves the norm and the denominator, while strictly reducing $|V|>1$.
Strong induction reaches $|V|=1$, where $|y_*|=1$.

Since $M\mid\overline Z$, reduction modulo $q^3$ and $q-1$ now gives
\[
 q^3\mid Y,\qquad q-1\mid B_t.
\]
No sign of $B_t$ has been used. The equation $A_1=0$ also gives an
integral main root $D_*/U$, whose absolute value can be classified.
Writing $\chi_A,\psi_A$ for the usual Pell sequences, positive
classification supplies indices $m_1,p\ge1$ such that
\[
 |\tau_*|=\chi_{2XY^2+1}(m_1),\qquad
 k=2\psi_{2XY^2+1}(m_1),\qquad c=\psi_{a+2}(p).
\]
Because $c\ge3$, $p\ge2$. The case $m_1=1$ would give
$k=2$, $\eta=\zeta=1$, $c=2Y+1$, contradicting
$c\ge2(a+2)\ge4Y+4$. Hence
\begin{equation}\label{eq:offset-early-c}
 c>kY\ge4(2XY^2+1)Y>8XY^3\ge8q^9.
\end{equation}
This is established before $C>0$, before the source bounds on $R$, and
before interpreting any recovered transport gap.

\subsection{The affine offset eliminates the exceptional-sign case}

Define the proof-only bound $\widehat R=q^3(q+F+\alpha)$.
The fixed-mask estimates give unconditionally
\[
 |R|<\widehat R.
\]
Put $h=F+\alpha\ge2$. Since $q\ge16$,
\[
 q^3(q+h)\le q^3(q/2+1)h<q^4h<ch
\]
by \eqref{eq:offset-early-c}. Thus
\[
 cT_{\rm old}-R=cT+ch-R>cT\ge c,
 \qquad V>D_fc-c>1.
\]
The new auxiliary coordinate is therefore positive at the required point
of the proof. This is what permits deletion of the predecessor's
$4qD_f$ exceptional-sign weight.

The same bound $ch>\widehat R$ also gives
\[
 0<cT_{\rm old}-R<c(T+2h),\qquad 0<V<D_fc(T+2h).
\]
Moreover,
\[
 \delta-(\Delta-1)D_f=\Delta(I^2-1)>0,
 \qquad c<k(Y+1),\qquad \Delta-1>2(Y+1).
\]
The last inequality follows from $a\ge2Y$ and $Y\ge1$. Therefore
\[
 0<y-V=\frac{V^2-1}{\delta(y+V)}<\frac{V}{\delta}
 <\frac{c(T+2h)}{\Delta-1}<\frac{k(T+2h)}2<k(T+h)=L.
\]
Also $f<C_f=(a+2)I$, since $C_f^2-D_f=I^2-1>0$. We conclude
\begin{equation}\label{eq:offset-small-L}
 L+C_f+V-f-y>0.
\end{equation}

Let $s=Y/q^3>0$ and $b=B_t/(q-1)\in\mathbb Z$. The recovered norm
identity reads
\[
 \tau_*+f_*+y_*=\ell+k+n+L+V+C_f+s-Wb.
\]
If $b\le0$, its right side is strictly larger than
$|\tau_*|+|f_*|+|y_*|$, because $|\tau_*|<\ell+k$ and
\eqref{eq:offset-small-L} holds. Hence $B_t>0$, which forces $C>0$.
Only now do we invoke the native source bounds $3q<R<q^4-2$.

Set
\[
 f_{\rm old}=|f_*|,\quad i_{\rm old}=i,\quad u=|y_*|-V,\quad
 r=n+L+V+C_f-f_*-y_*.
\]
All are positive integers, with $r>n$. They satisfy the weighted source70
carrier with weight $h_*=L$, its normalized strong norm, its auxiliary
norm, and its adaptive main norm. That inherited interface is sound for
every positive integer $h_*$. Its native/compiler and syntactic-refutation
calibration therefore supplies a finite ZFC refutation.

\subsection{Completeness and the exact cost of the offset}

A finite refutation gives the canonical compact source, with
$R\equiv3\pmod4$, $R>3q$, $F+\alpha<q$, and the earlier bound on $c$.
Choose its normalized strong pair $f,I$, $I=ic^2$, and write
$M_*=\Delta I$. The inherited positive auxiliary reconstruction uses
\[
 V_* =\frac{\chi_{M_*}(R)}{M_*},\qquad
 y_* =\psi_{M_*}(R),\qquad
 T_{\rm old}=\frac{V_*+c+Rf^2}{cf^2}.
\]
It proves the last quantity integral and $y_*>V_*>0$. The stronger bound
needed for the new offset is elementary. Since $R\ge7$,
\[
 V_*\ge M_*^{R-1}\ge M_*^6,\qquad f^2<M_*^2,\qquad M_*\ge c^2,
\]
and hence
\[
 T_{\rm old}>\frac{V_*}{cf^2}>\frac{M_*^4}{c}\ge c^7>q>F+\alpha.
\]
Thus $T=T_{\rm old}-F-\alpha>0$ is integral. No new cofinality assumption
or congruence on the strong Pell pair has been introduced.

The same growth also proves $T_{\rm old}>C_f$: because $I\ge c^2$,
\[
 T_{\rm old}>\frac{M_*^4}{c}=\Delta^4I\frac{I^3}{c}>(a+2)I=C_f.
\]
After these choices compute $L,W$. We have $L=kT_{\rm old}>C_f$ and
$L>k$. The canonical transport has $t_0=B_t/(q-1)\ge q$, $s<k$, and
$\tau>\ell$, so
\[
 r=Wt_0+\tau-\ell-k-s>qW-2k>2L.
\]
Consequently
\[
 n=r-L-V_*-C_f+f+y_*=r-L-C_f+f+(y_*-V_*)>0.
\]
This makes $\overline Z=M(\tau+f+y_*)$ and $N_8=0$; the inherited
adaptive carrier supplies positive $\gamma$ and $A_1=0$. Completeness
is proved on canonical witnesses, which is exactly what an existence
equivalence requires.

Both residuals are squares modulo four. Thus the factorization
\[
 1+P_{11}=(1+N_8)(1+A_1)
\]
shows $P_{11}=0$ if and only if $N_8=A_1=0$: a negative unit would
require a residual $-2$, impossible modulo four.

\subsection{Exact degree and verification}

The degree audit uses an exact universal double-center identity. For
independent variables $z,a,b,c,e,f$, set
\[
 x=z+b+c,\quad B=b^2-e,\quad C=c^2-f,\quad
 N_c^\circ=(x^2-a-B-C)^2-4(aB+aC+BC),
\]
\[
 N^\circ=(N_c^\circ)^2-64aBCx^2.
\]
Give $(z,a,b,c,e,f)$ the weights $(8,16,11,22,18,26)$.
Exact expansion, recorded in a $229$-term certificate, gives
\[
 [N_c^\circ]_{63}=8bc^2z,\qquad
 [N^\circ]_{126}=64b^2c^4(z^2-a).
\]
Substitute
\[
 z=\overline Z-MV-MC_f,\quad a=A,\quad b=MC_f,\quad c=MV,
\]
\[
 e=M^2(I^2-1),\qquad f=-\frac{M^2(V^2-1)}\delta.
\]
These are exact because $C_f^2-D_f=I^2-1$ and
$(MV)^2-C_{\rm num}/\delta=-M^2(V^2-1)/\delta$. Multiplication by
$\delta^2$ and $\delta^4$ gives the actual inner and outer norm.
The power of $f$ is at most four, so the cleared identity is polynomial.
Since $\deg\delta=18$, the exact residual degrees are
\[
 \deg N_c=99,\qquad \deg N_8=198,\qquad\deg A_1=13,
 \qquad\deg P_{11}=211.
\]

For an attained-degree certificate on the actual fixed tuple, put
\begin{align*}
 q_0&=3\beta j,\qquad k_0=\eta+\zeta,\qquad I_0=ik_0^2Y^2,\\
 D_{f,0}&=X^2Y^2I_0^2,\qquad Q_0=X^4Y^4I_0^2,\qquad C_{f,0}=XYI_0,\\
 L_0&=k_0(T+F+\alpha),\qquad\ell_0=Xk_0Y^2,\qquad R_0=q_0^3(q_0-F),\\
 B_{t,0}&=(K_0q_0+X)(q_0-F-\alpha)-q_0F.
\end{align*}
Then
\begin{align*}
 [N_8]_{198}&=64Q_0^4q_0^{32}D_{f,0}^4R_0^4C_{f,0}^2
                 L_0B_{t,0}(L_0B_{t,0}-2\ell_0),\\
 [A_1]_{13}&=32q_0^2k_0X^2Y^4(4q_0^2Y\gamma+R_0).
\end{align*}
The coefficient of
\[
 j^{48}i^{18}\eta^{39}X^{30}Y^{67}T^2F^4\alpha^2\gamma
\]
is $2^{13}(3\beta)^{48}>0$, independently of $K_0,M_C,M_F$.
The literal graph has $64$ multiplications, $31$ additions, and $17$
subtractions. Scalar multiplication and squaring are charged; fixed
numeral leaves and register reuse are free. The two-residual system has
$109$ operations, and the single equation has $112$.

The separate exact audit checks every named polynomial contract, the exact
double-center cancellation, the eight-conjugate identity, radical-extraction
identities, the attained leading coefficient, liveness, and operation
counts. The complete compact44 payload is verified and independently
reconstructed from its finite recipe. Normal and optimized Python runs
pass with explicit checks.

Combining soundness and completeness with the named inherited source
proves the ordinary mathematical consistency equivalence on the stated
positive domain. It does not lower the existing nine-positive variable
count, prove global minimality, or provide a new end-to-end formal
calibration. The new content is the offset, the resulting early sign
argument, its canonical completeness proof, and the degree-$211$ endpoint.

